
\begin{tikzpicture}[
    >=Latex,
    font=\small,
    block/.style={
        draw,
        rounded corners=2pt,
        align=center,
        minimum width=2.5cm,
        minimum height=0.9cm,
        inner sep=5pt
    },
    smallblock/.style={
        draw,
        rounded corners=2pt,
        align=center,
        minimum width=2.3cm,
        minimum height=0.8cm,
        inner sep=5pt
    },
    arrow/.style={->, thick}
]

% Recepción analógica
\node[block] (mic1) at (0,2.0) {Micrófono 1};
\node[block] (mic2) at (0,0.5) {Micrófono 2};

\node[block] (xadc) at (3.4,1.25)
    {XADC\\Conversión A/D};

\node[block] (zynq) at (6.9,1.25)
    {Zynq-7000\\PYNQ-Z2};

\node[smallblock] (dma) at (10.2,2.15)
    {DMA\\Transferencia};

\node[smallblock] (processing) at (10.2,0.35)
    {Procesamiento\\digital};

\node[block] (software) at (13.6,2.15)
    {Software\\de control};

\node[smallblock] (analysis) at (13.6,0.35)
    {Análisis\\de señales};

% Conexiones de recepción
\draw[arrow] (mic1.east) -- (xadc.west);
\draw[arrow] (mic2.east) -- (xadc.west);

\draw[arrow] (xadc.east) -- (zynq.west);
\draw[arrow] (zynq.east) -- (dma.west);
\draw[arrow] (zynq.east) -- (processing.west);

\draw[arrow] (dma.east) -- (software.west);
\draw[arrow] (processing.east) -- (software.south west);
\draw[arrow] (software.south) -- (analysis.north);

% Ruta de excitación desde la tarjeta
\node[smallblock] (pulse) at (3.4,-1.25)
    {Pulso digital\\de salida AR2};

\node[smallblock] (driver) at (7.1,-1.25)
    {Etapa de excitación\\Transistor 2N2222A};

\node[block] (buzzer) at (10.6,-1.25)
    {Buzzer activo};

\draw[arrow] (pulse.east) -- (driver.west);
\draw[arrow] (driver.east) -- (buzzer.west);

% Relación de la salida digital con la tarjeta
\draw[arrow, dashed] (zynq.south) -- ++(0,-0.55)
    -| (pulse.north);

\end{tikzpicture}
